% Part: history
% Chapter: biographies
% Section: emmy-noether

\documentclass[../../../include/open-logic-section]{subfiles}

\begin{document}

\olfileid{his}{bio}{noe}

\olsection{এমি নোয়েটার}

\olphoto{noether-emmy}{এমি নোয়েটার}

এমি নোয়েটারের (উচ্চারণ: \textsc{নের}-টার)
জন্ম ১৮৮২ সালের ২৩ মার্চ জার্মানির এরলাঙ্গেনে,
একটি বিদ্বৎ ও উচ্চ-মধ্যবিত্ত পরিবারে। তাঁকে
“আধুনিক বীজগণিতের জননী” বলা হয়। নারীশিক্ষার
পথে বড় বাধা সত্ত্বেও তিনি গণিত ও পদার্থবিদ্যা,
উভয় ক্ষেত্রেই যুগান্তকারী অবদান রাখেন। সেই সময়ে
জার্মানিতে মেয়েদের কলাশিক্ষার জন্য প্রস্তুত করা
হতো; কলেজের প্রস্তুতিমূলক বিদ্যালয়ে তাদের
ভর্তির অনুমতি ছিল না। তবু গ্যোটিঙেন ও এরলাঙ্গেন
বিশ্ববিদ্যালয়ে (শেষোক্তটিতে তাঁর বাবা গণিতের
অধ্যাপক ছিলেন) অনুমতি নিয়ে ক্লাসে বসার পর,
১৯০৪ সালে এরলাঙ্গেনের নিয়ম বদলে ছাত্রীদের
ভর্তির সুযোগ হলে নোয়েটার সেখানে শিক্ষার্থী
হিসেবে ভর্তি হতে পারেন। ১৯০৭ সালে তিনি
গণিতে ডক্টরেট লাভ করেন।

যোগ্যতা থাকা সত্ত্বেও কর্মজীবনে নোয়েটারকে অনেক
বাধার মুখে পড়তে হয়। ১৯০৮--১৯১৫ সালে তিনি
এরলাঙ্গেনে বিনা বেতনে পড়ান। এই সময়ে তৎকালীন
বিশ্বের অন্যতম প্রধান গণিতবিদ ডাভিড হিলবার্টের
নজরে পড়লে তিনি নোয়েটারকে গ্যোটিঙেনে আমন্ত্রণ
জানান। কিন্তু মহিলাদের অধ্যাপক হওয়া নিষিদ্ধ
ছিল; ফলে সেখানেও তিনি হিলবার্টের নামে, বিনা
বেতনে বক্তৃতা দিতে পারতেন মাত্র। এই সময়েই
তিনি বর্তমানে নোয়েটারের উপপাদ্য নামে পরিচিত
ফলটি প্রমাণ করেন, যা আজও তাত্ত্বিক পদার্থবিদ্যায়
ব্যবহৃত হয়। অবশেষে ১৯১৯ সালে নোয়েটার
পাঠদানের অধিকার পান। বিশ্ববিদ্যালয়ের
সহকর্মীদের অব্যাহত আপত্তির জবাবে হিলবার্ট
নাকি বলেছিলেন, “ভদ্রমহোদয়েরা, অনুষদ-পরিষদ
কোনো স্নানাগার নয়।”

১৯২০-এর দশকের শেষ দিকে নোয়েটার বিমূর্ত
বীজগণিতে মনোনিবেশ করেন এবং তাঁর অবদান
ক্ষেত্রটিকে বদলে দেয়। প্রমাণে তিনি প্রায়ই
তথাকথিত \emph{ঊর্ধ্বগামী শৃঙ্খল শর্ত} ব্যবহার
করতেন। এই শর্ত বলে, নির্দিষ্ট ধরনের সেটের
কোনো অসীম কঠোরভাবে বর্ধমান শৃঙ্খল নেই।
যেমন, বর্তমানে নোয়েটারীয় বলয় নামে পরিচিত
কিছু বীজগাণিতিক কাঠামোতে আইডিয়ালের
$I_1 \subsetneq I_2 \subsetneq \dots$ আকারের
কোনো অসীম অনুক্রম নেই। শর্তটি যেকোনো আংশিক
ক্রমে সাধারণীকরণ করা যায় (বীজগণিতে
সেট-অন্তর্ভুক্তি অনুযায়ী ক্রমবদ্ধ আইডিয়ালগুলির
বিশেষ ক্ষেত্রটি বিবেচিত হয়)। এর দ্বৈত
\emph{নিম্নগামী শৃঙ্খল শর্ত}ও বিবেচনা করা যায়:
আংশিক ক্রমে প্রতিটি কঠোরভাবে \emph{হ্রাসমান}
অনুক্রম শেষ পর্যন্ত থামে। কোনো আংশিক ক্রম
নিম্নগামী শৃঙ্খল শর্ত পূরণ করলে, যেমনভাবে
~$\Nat$-এ $<$ ক্রম বরাবর আরোহ প্রয়োগ করা
যায়, তেমনি ওই ক্রম বরাবরও আরোহ প্রয়োগ
করা যায়। এমন ক্রমকে \emph{সুপ্রতিষ্ঠিত}
বা \emph{নোয়েটারীয়} বলা হয়; সংশ্লিষ্ট
প্রমাণনীতির নাম \emph{নোয়েটারীয় আরোহ}।
% BN-SRC-848: অতিরিক্ত ব্যাখ্যা/অনিশ্চয়তা review-ledger-এ; reader-এ প্রয়োগ নয়।

নোয়েটার ছিলেন ইহুদি। ১৯৩৩ সালে নাৎসিরা
ক্ষমতায় এলে তাঁকে তাঁর পদ থেকে বরখাস্ত করা
হয়। সৌভাগ্যক্রমে, তিনি যুক্তরাষ্ট্রে চলে যেতে
পারেন এবং পেনসিলভানিয়ার ব্রিন মাওরে একটি
অস্থায়ী পদ পান। সেখানে থাকার সময়ে তিনি
প্রিন্সটনের ইনস্টিটিউট ফর অ্যাডভান্সড স্টাডিতেও
বক্তৃতা দিতেন। প্রিন্সটন বিশ্ববিদ্যালয়কে
তাঁর কাছে নারীদের প্রতি অনাতিথেয় মনে হয়েছিল
\citep[81]{Dick1981}। ১৯৩৫ সালে জরায়ুর
একটি টিউমার অপসারণের জন্য তাঁর অস্ত্রোপচার
হয়। অস্ত্রোপচারের ফলে সংক্রমণ ঘটলে তাঁর মৃত্যু
হয়; তাঁকে ব্রিন মাওরেই সমাহিত করা হয়।
% BN-SRC-849 TARGET-CORRECTION-BEGIN
\par\noindent\textbf{পাঠ-সংশোধন।} প্রিন্সটনে তাঁর অতিথি-পাঠদানের প্রতিষ্ঠানটি
ইনস্টিটিউট ফর অ্যাডভান্সড স্টাডি; এটি প্রিন্সটন
বিশ্ববিদ্যালয় থেকে পৃথক। দেখো
\href{https://celebratio.org/Noether_E/article/111/}{ভাইলের স্মৃতিবক্তৃতা}।
% BN-SRC-849 TARGET-CORRECTION-END

\begin{reading}
নোয়েটারের জীবনী জানতে \citet{Dick1981}
দেখো। পেরিমিটার ইনস্টিটিউট ফর থিওরেটিক্যাল
ফিজিক্সের ওয়েবসাইটে তাঁর জীবন ও প্রভাব নিয়ে
প্রতিষ্ঠানটির বক্তৃতাগুলি পাওয়া যায়
\citep{Perimeter2015}। পড়তে পড়তে ক্লান্ত
হলে নোয়েটারের জীবন ও প্রভাব নিয়ে
\emph{Stuff You Missed in History Class}
পডকাস্টটি শুনতে পারো \citep{Frey2015}।
নোয়েটারের সংকলিত রচনা মূল জার্মান ভাষায়
পাওয়া যায় \citep{Noether1983}।
\end{reading}

\end{document}
